% Generated draft: requires mathematical and source review.
\documentclass[11pt]{article}
\usepackage[margin=25mm]{geometry}
\usepackage{fontspec}
\setmainfont{Latin Modern Roman}
\usepackage{amsmath,amssymb,mathtools,unicode-math}
\setmathfont{Latin Modern Math}
\usepackage{enumitem,needspace,xurl,hyperref}
\setlength{\parindent}{0pt}
\setlength{\parskip}{5pt}
\setlength{\emergencystretch}{3em}
\newcommand{\N}{\mathbb{N}}
\newcommand{\Z}{\mathbb{Z}}
\newcommand{\Q}{\mathbb{Q}}
\newcommand{\R}{\mathbb{R}}
\newcommand{\C}{\mathbb{C}}
\newcommand{\F}{\mathbb{F}}
\newcommand{\A}{\mathbb{A}}
\newcommand{\PP}{\mathbb{P}}
\newcommand{\E}{\mathbb{E}}
\newcommand{\eps}{\varepsilon}
\newcommand{\dd}{\,\mathrm{d}}
\begin{document}
\section*{UnsolvedMath ID: 2307036}
\textbf{Source problem code:} catalogue entry
\subsection{Definitions and mathematical statement}
Let $\mathbb D=\{z\in\C:|z|<1\}$, let $m$ be plane Lebesgue measure, and let $E\subseteq\mathbb D$ be measurable. The source defines the two-dimensional Hilbert transform of $\mathbf1_E$ by
\[
 f_E(z)=-\frac1\pi\iint_E\frac{\dd m(w)}{(w-z)^2},\qquad z\in\C\setminus E.
\]
For arbitrary measurable $E$, points outside $E$ can still be arbitrarily close to its mass, so the literal integral needs an interpretation. The standard singular-integral reading is the almost-everywhere principal value
\[
 f_E(z)=\lim_{\epsilon\downarrow0}-\frac1\pi
 \iint_{E\setminus B(z,\epsilon)}\frac{\dd m(w)}{(w-z)^2}.
\]
Away from the closure of $E$ the ordinary integral suffices. Source review should confirm the intended almost-everywhere version; the estimate concerns its integral over $\mathbb D\setminus E$.

The source reports constants $a,b$ such that
\[
 \iint_{\mathbb D\setminus E}|f_E|\,\dd m
 \leq a\,m(E)\log\frac\pi{m(E)}+b\,m(E)
\]
for every measurable $E\subseteq\mathbb D$, by Calderón--Zygmund theory. Determine the smallest admissible leading coefficient $a$, where the additive constant $b$ may depend on $a$ but must be independent of $E$. The conjectured sharp endpoint is
\[
 \exists b\in\R\ \forall E\subseteq\mathbb D\text{ measurable},\quad
 \iint_{\mathbb D\setminus E}|f_E|\,\dd m
 \leq m(E)\log\frac\pi{m(E)}+b\,m(E).
\]
At $m(E)=0$, use the continuous convention $m(E)\log(\pi/m(E))=0$. The source says small centered disks show that a coefficient below one is impossible. Proving only inequalities with every $a>1$ would not by itself establish the displayed endpoint with $a=1$.

The source calls this equivalent to Gehring--Reich's area-distribution Problem 7.9, whose statement is not supplied. It also reports a sharp one-dimensional analogue for $E\subset[-1,1]$ and the transform $-\pi^{-1}\int_E\dd t/(x-t)$, by subordination or a Stein--Weiss theorem; no exact analogue inequality is included.
\subsection{Short English statement}
Find the sharp leading logarithmic constant in the integral estimate for the planar Hilbert transform of an indicator set. The conjectured coefficient is one with a single additive constant working for every measurable subset of the disk. The registry labels this solved, while its research summary says OPEN-TRIAGE and no definitive resolution located after Hayman's 2018 update. No resolving citation is supplied, so that status conflict is retained rather than silently resolved.
\subsection{Sources}
\begin{itemize}
\item UnsolvedMath permanent ID 2307036, source code catalogue entry, ``Research Problems in Function Theory --- Problem 7.36.'' Registry curation and metadata: Ulam AI and UnsolvedMath Contributors, CC BY 4.0; updated 2026-08-06.
\item Hayman, \emph{Research Problems in Function Theory} (2018 source update), Problem 7.36, attributed to A. Baernstein, comparing Gehring--Reich's Problem 7.9: \url{https://arxiv.org/abs/1809.07200}. Supplied source-TeX extraction accessed 2026-07-29. The metadata retains the update for status triage and marks status/release and rights review NEEDS\_REVIEW.
\item Calderón--Zygmund theory, Gehring--Reich, and Stein--Weiss are named in the source without full bibliographic entries. No missing theorem, one-dimensional formula beyond the supplied transform, or resolving reference is inferred.
\end{itemize}
\end{document}
